% Four Measurement Failures in a Small-Scale Language Model Experiment
%
% Self-contained so it compiles with a standard TeX install. To retarget at an
% official venue class, replace the geometry/titling block with the venue's
% \usepackage line; the body needs no other changes.

\documentclass[10pt,twocolumn]{article}

\usepackage[letterpaper,margin=0.75in,columnsep=0.25in]{geometry}
\usepackage{times}
\usepackage{booktabs}
\usepackage{amsmath}
\usepackage{microtype}
\usepackage[hidelinks]{hyperref}
\usepackage{titlesec}
\usepackage{titling}

\titleformat{\section}{\normalfont\normalsize\bfseries}{\thesection}{0.5em}{}
\titlespacing*{\section}{0pt}{0.9em}{0.3em}
\setlength{\parindent}{1em}
\setlength{\droptitle}{-4em}
\renewcommand{\baselinestretch}{0.98}
\emergencystretch=1.5em

\title{\bfseries Four Measurement Failures in a Small-Scale\\
Language Model Experiment}
\author{}
\date{}

\begin{document}
\maketitle
\vspace{-5em}

\begin{abstract}
\noindent
A common argument for filtering fact-heavy pages out of language model
pretraining data is that memorising facts consumes parameters the model could
otherwise spend on reasoning. We set out to test this claim directly at small
scale, and over four rounds of experiments we obtained four null results. Each was
uninterpretable, because in every round a different component of our
measurement pipeline had failed in a way our existing checks did not detect: a decoder that scored a working benchmark at 4.7\% when its true score
was 93.8\%, a data pipeline that replaced most of one training stream with
another, a control condition that removed less learnable signal than the
treatment it was matched against, and a knowledge probe that could not
distinguish a fact the model had memorised from one we invented. In this paper
we report each failure, quantify it, and give a cheap test that detects it. We
also report the one measurement that survived. Under teacher forcing, loss at
fact positions is the same for people the model trained on as for people it
never saw, differing by $-0.0010 \pm 0.0132$ nats. The model has fitted the
biography format and the frequencies of the value pools, and we find no
evidence that it stored anything specific to individual people. We first
reached that conclusion with a broken probe, and only later built an instrument
that can detect storage when storage is present.
\end{abstract}

\section{Introduction}

Language models are increasingly trained on curated rather than scraped data,
and one recurring curation decision is how much factual content to include.
Several authors have suggested that arbitrary facts are a poor use of a small
model's parameters, on the grounds that memorising them consumes capacity that
would otherwise support reasoning. If this is true, then removing or hiding
factual content during training should improve reasoning performance at fixed
model size, and the effect should grow with the amount of factual content
removed.

We attempted to measure this effect directly. Working at 40M parameters, where
a full training run costs about a day on one GPU, we ran four rounds of the
experiment. Each returned a null result: no detectable difference between the
treatment condition and its control. Because null effects are plausible at this
scale, we did not initially treat repeated nulls as evidence that the pipeline
itself had failed.

Each null was in fact uninterpretable, for a different reason each time. This
paper reports those four failures rather than the hypothesis, which we still
regard as untested. For each we give the magnitude and a test that would have
caught it; none of the tests requires a GPU and the most expensive takes fifteen
minutes on a laptop. Three of the failures are properties of tooling that many
groups write for themselves. The fourth is of a kind the usual defences do not
catch: that measurement had a control attached, the control passed, and the
measurement was still worthless, because nothing had shown it could return a
positive under any circumstances.

Section~\ref{sec:bg} describes the model, the data, and the metrics.
Sections~\ref{sec:pad} through~\ref{sec:store} report the four failures in the
order we found them. Section~\ref{sec:disc} discusses what they have in common,
and Section~\ref{sec:lim} states the limits of the evidence.

\section{Background and Setup}
\label{sec:bg}

\paragraph{Model and data.}
We trained a 40.6M-parameter decoder-only transformer from scratch. The
training corpus is synthetic, which lets us control exactly how much factual
content the model sees. It interleaves four streams of text. The first is short
biographies of invented people, each stating six attributes such as a birth
date and an employer, with values drawn from fixed pools. These are the
``facts'' whose effect we wanted to measure. The second is modular-arithmetic
word problems accompanied by step-by-step worked solutions, following the
design of Ye et al.~\cite{igsm}. The third is chains of logical deduction. The
fourth is ordinary natural text, included as filler so that the total size of
the corpus can be held fixed while the proportions of the other three vary.

The main run described here trained for 31{,}280 optimisation steps on 16.4
billion tokens, and its corpus contains 1{,}531{,}800 invented people, each
appearing 100 times.

\paragraph{Reporting accuracy.}
All accuracies below are \emph{closed-book exact match}: the model sees a
problem with no reference material and its answer counts as correct only if it
matches the target string exactly. We compare against the \emph{majority-class
baseline}, the accuracy obtained by always giving whichever answer is most
common in the test set. Our arithmetic answers are remainders after division by
a small integer, so some occur far more often than others, and comparing
against uniform chance would flatter the model.

\paragraph{Measuring loss.}
Losses are in \emph{nats}, the natural-logarithm form of cross-entropy. A loss
of $\ln k$ nats means the model is as uncertain as if choosing uniformly among
$k$ options, so lower is better and a change of 0.69 nats is a factor of two in
the probability assigned to the correct token.

\paragraph{Storage capacity.}
Allen-Zhu and Li~\cite{az} estimate how much factual knowledge a transformer can
hold, finding a ceiling near two bits per parameter when each fact appears about
a thousand times in training, and near one bit at a hundred appearances. Our
model sees each fact 100 times, and its corpus would need about two bits per
parameter to store completely, so the corpus asks for roughly twice what the
model can hold. Our operating point coincides with one of their reported
measurements rather than falling between them, so this ratio does not depend on
interpolating their curve, and their hundred-appearance figure would have to
double before our corpus stopped exceeding capacity.

\section{A Decoder That Attended to Its Own Padding}
\label{sec:pad}

To evaluate the model we generate answers for many problems at once. Prompts
have different lengths, so batching them requires padding to a common length.
Our generator padded on the left, used the end-of-text token as the padding
symbol, and supplied no \emph{attention mask} --- the array telling a
transformer which positions to ignore. Without one the model attends to padding
as though it were content. A comment in our code justified this on the grounds
that the pads were ``ordinary end-of-text tokens the model has seen as
separators,'' which is exactly the problem. During training that token marks a
boundary between documents, so a prompt preceded by a run of them looks like the
opening of a new document, and the model does what it was trained to do there:
it writes one. The output is fluent, correctly formatted, and parses cleanly as
an answer, which is why the failure survived our output checks. The model was
answering a question it had invented.

Table~\ref{tab:pad} shows the effect on a single checkpoint, evaluated three
ways on the same 64 held-out problems.

\begin{table}[h]
\centering\small
\begin{tabular}{lr}
\toprule
decoding method & exact match \\
\midrule
batched, padded to batch maximum & 3/64 \\
one prompt at a time & 60/64 \\
batched, after the fix & 60/64 \\
\bottomrule
\end{tabular}
\caption{The same weights on the same items, decoded three ways.}
\label{tab:pad}
\end{table}

A replication at a different random seed gave 0/64, 63/64 and 63/64, and every
repaired batched generation was character-for-character identical to its
unbatched counterpart. By varying the amount of padding on a fixed prompt we
found that corruption sets in somewhere between 16 and 24 padding tokens.

The defect also explains two anomalies we had attributed to the model itself.
Measured accuracy had been \emph{rising} with the number of arithmetic steps a
problem required, the opposite of the expected pattern, which we had read as a
model emitting a near-constant answer. In fact padding fills the gap between a
prompt and the longest prompt in its batch, and harder problems have longer
prompts, so the easiest problems received the most padding and were damaged
most. Separately, 47\% of generations had fallen outside the set of
syntactically valid answers, which we had taken as evidence that the model
could not produce a well-formed response.

The general fix is to supply an attention mask, which costs nothing. The test
is one comparison: decode a prompt on its own, decode it again inside a batch
of mixed-length prompts, and require identical output.

With the decoder repaired the benchmark works. Table~\ref{tab:ops} reports the
finished run on 1{,}500 problems from the difficulty range it trained on and 750
harder ones. Aggregate accuracy in the trained range is 96.0\%, against a
majority-class baseline of 7.5\%, and the deduction task reaches 0.950 against a
baseline of exactly 0.500. One to three operations are saturated and no longer
informative; four still has room, and the sharpest remaining signal is the fall
from five operations to six.

\begin{table}[h]
\centering\small
\begin{tabular}{clrr}
\toprule
operations & in training range & $n$ & exact match \\
\midrule
1 & yes & 399 & 99.7\% \\
2 & yes & 384 & 99.2\% \\
3 & yes & 375 & 95.7\% \\
4 & yes & 342 & 88.3\% \\
5 & no  & 183 & 41.5\% \\
6 & no  & 194 & 7.2\% \\
7 & no  & 190 & 5.3\% \\
8 & no  & 183 & 4.9\% \\
\bottomrule
\end{tabular}
\caption{Final checkpoint, by problem difficulty.}
\label{tab:ops}
\end{table}

\section{A Control Condition With Incompatible Criteria}
\label{sec:ctrl}

The experiment compares two training conditions. In the treatment, the loss is
not computed on the tokens that spell out a fact's value, so the model is not
asked to predict them. In the control, the same number of tokens is excluded
from the loss, but they are ordinary text rather than fact values. The point of
the control is to separate the effect of \emph{which} tokens were excluded from
the effect of simply excluding some.

For the comparison to isolate the first effect, the excluded spans in the two
conditions should match on everything except their content. Our match
constrained four properties: how many tokens are excluded, how long each
contiguous excluded span is, where in the document those spans sit, and how
difficult the excluded tokens are to predict. The last is the important one,
since excluding easy tokens removes less learnable signal than excluding hard
ones.

Measured across the corpus, the match failed. Treatment spans averaged 1.9598
nats and control spans 1.2712, a relative shortfall of 35.1\% against a
tolerance of 20\% that we had fixed before running the experiment. Our first
interpretation was that the corpus simply contains no ordinary text difficult
enough to match fact values. Table~\ref{tab:ctrl} shows that this is wrong, and
the way it is wrong is more useful than the original observation. Relaxing the
matching criteria one at a time shows that sufficiently difficult tokens do
exist, but only when the requirement that they form contiguous spans is
dropped.

\begin{table}[h]
\centering\small
\begin{tabular}{lrr}
\toprule
control construction & mean nats & shortfall \\
\midrule
as shipped: length + position & 1.2712 & 35.1\% \\
best possible, length only & 1.5667 & 20.0\% \\
scattered, difficulty-matched & 1.7677 & 9.8\% \\
scattered, hardest available & 1.8644 & 4.9\% \\
\bottomrule
\end{tabular}
\caption{Pricing each matching criterion separately, over 800 documents.}
\label{tab:ctrl}
\end{table}

The reason this generalises is that facts and ordinary prose distribute their
surprising tokens differently. A fact value is a contiguous run of hard-to-predict
tokens; comparably hard tokens in ordinary text occur singly and scattered. A
control can therefore reproduce the shape of what the treatment removes, or the
quantity of learnable signal it removes, but not both at once. We do not think
this is specific to our corpus, and we suggest that work using matched controls
of this kind state which of the two properties it chose to preserve. We had not
realised we were choosing.

\section{A Data Stream That Silently Became Another}
\label{sec:mix}

Our corpus builder mixes the four streams according to requested proportions.
The fact stream is finite: with a fixed number of invented people, each
appearing a fixed number of times, it can emit only so many documents. When the
requested proportion demanded more than that, the builder made up the shortfall
from the filler text so that the total token count remained exact. Holding the
token count fixed is a requirement of the design, because every experimental
condition must train for the same number of steps.

This reallocation is the correct behaviour and it produced no warning.
Table~\ref{tab:mix} shows what one pilot corpus was asked for and what it
actually contained.

\begin{table}[h]
\centering\small
\begin{tabular}{lrr}
\toprule
stream & requested & realised \\
\midrule
facts & 50\% & 3.75\% \\
filler text & 10\% & 56.25\% \\
\bottomrule
\end{tabular}
\caption{A corpus built for an experiment about fact load.}
\label{tab:mix}
\end{table}

The pilot that produced our first decision to stop the experiment had therefore
trained on a corpus that was 3.75\% facts. It passed every integrity check we
had, because each check tested the corpus for internal consistency --- that the
token counts add up, that the masks align with the text --- rather than testing
it against its own specification. The realised proportions were written to a
metadata field that nobody read.

The remedy is to record the requested proportions next to the realised ones and
to fail the build on any discrepancy. The requested values have to be stored
separately from the working budget, because the reallocation overwrites the
budget it draws from.

\section{A Knowledge Probe That Could Not Detect Knowledge}
\label{sec:store}

The hypothesis under test assumes the model stores the facts it is shown, so we
checked. Loss at fact positions in training documents is 1.9732 nats, against
25.03 on an earlier corpus where facts appeared too rarely to be learned. That
improvement is real but smaller than it looks: 1.9732 nats per token works out
to 53.08 bits per person, against 52.96 bits for a predictor that knows only
how often each value occurs in the corpus. The model has reached the value-pool
frequencies and no further.

To ask whether it had nonetheless stored the individual associations, we
queried it with a phrasing that never appears in training and converted its
answers into \emph{recoverable bits}: how much better it predicts the true
value than that frequency baseline would. It scored well below zero, and so did
a set of invented people it had never seen. We wrote this up as a model that
had stored nothing retrievable.

That description is correct, as we show below, but the evidence we had for it
was worthless. We discovered this by asking the probe for a result that had to
come back positive. If a probe can detect stored knowledge at all, then it must
separate people the model trained on from people we invented when the query
uses a phrasing the model saw during training. Table~\ref{tab:probe} shows that
it does not.

\begin{table}[h]
\centering\small
\begin{tabular}{lrrr}
\toprule
query phrasing & trained & invented & diff. \\
\midrule
absent from training & $-111.93$ & $-112.10$ & $+0.17$ \\
present in training & $-67.83$ & $-67.65$ & $-0.18$ \\
\bottomrule
\end{tabular}
\caption{Recoverable bits per person. The probe cannot tell the two groups
apart even where it must.}
\label{tab:probe}
\end{table}

The arithmetic above points the same way. Training loss puts the model within
0.12 bits of the frequency baseline, whereas the probe reports it as 68 bits
per person worse. Two measurements of the same quantity disagreed by a wide
margin and we had quoted the one that suited our expectations.

\paragraph{A measurement that does not use a query.}
We therefore measured storage in the terms the training objective itself uses.
We take a training document, present it to the model as training did, and read
the loss only at the positions spelling out the facts. This is \emph{teacher
forcing}: the model receives the true preceding tokens rather than its own
generated ones, so the loss measures prediction given correct context. If the
model memorised a fact, its document must be cheaper at those positions,
because that is where knowing the fact pays. A model that learned only which
values are common gets no such discount.

\begin{table}[h]
\centering\small
\begin{tabular}{lrr}
\toprule
people & nats/token & docs \\
\midrule
trained, 100 appearances & $1.9876\!\pm\!0.0096$ & 400 \\
invented, never seen & $1.9866\!\pm\!0.0091$ & 400 \\
difference & $-0.0010\!\pm\!0.0132$ & $z\!=\!-0.08$ \\
\bottomrule
\end{tabular}
\caption{Standard errors are computed over documents. Each person contributes
two renderings, so clustering by person would widen them by up to a factor of
$\sqrt{2}$; the bounds we quote below allow for that.}
\label{tab:store}
\end{table}

\paragraph{Validating the replacement.}
This measurement is subject to the same objection as the probe it replaces: a
null result from an instrument that has never returned anything else is not
evidence. We therefore constructed a case in which the answer must be positive.
A small model trained for 250 steps on nothing but 48 biographies has certainly
memorised them. On that model the same measurement separates the two groups by
1.4484 nats ($z=+18.7$), and training it further widens the gap rather than
closing it, which establishes sensitivity to a large gap in this overfitting
regime.

In the memorising case both groups draw their values from the same pools, so
knowing which values are common cannot account for the separation. On the 40M
model the same measurement finds no such difference, and after widening the
interval for clustering it excludes any advantage larger than about 0.04 nats,
roughly a fortieth of the gap it resolves on the memorising model.

We are careful about what this licenses. The toy run shows the instrument can
detect entity-specific storage when that storage is forced into a small closed
set. It does not show that sparse or partial storage across 1.5M people would
appear as a teacher-forced gap of this size, and an average null can hide a
memorised subset. The comparison also changes two things at once, since
invented people differ from trained ones both in their attribute associations
and in whether their names were seen at all. The cleaner test holds the people
fixed and permutes their values among them, which we have not run. What we can
say is narrower: under the training objective, at fact positions, people the
model trained on are indistinguishable on average from people we invented,
which is what one would expect if prediction there were driven by template and
pool cues rather than by who the person is.

If the capacity estimate transfers to this setting, a load at twice that
estimate predicts storage of about half the facts, whereas we measure a mean
advantage indistinguishable from zero. Resolving that needs the
capacity-matched run described below, so we report it as a discrepancy rather
than as a result about capacity.

\section{Discussion}
\label{sec:disc}

The hypothesis we set out to test assumes that stored facts occupy capacity
reasoning could otherwise use. In our setting we can detect no such stored
facts, so there was nothing for the treatment to free and the experiment cannot
speak to the hypothesis in either direction.

Three of the four failures are ordinary engineering defects. We report them
because each had a cheap detector that we did not run, and because the tooling
they occurred in is the kind most groups write for themselves. We do not claim
they share a mechanism.

The fourth is different, and it is the one we think generalises. A negative control is the
usual defence for a measurement like our knowledge probe, and we had one: the
probe scored invented people no differently from real ones. But a broken probe
and an empty model give identical readings under that control, so it cannot
distinguish them. The control that separates them is the one that forces a
positive result. We suggest the requirement is general: before a null result is
allowed to bear on a hypothesis, the instrument that produced it should be shown
capable of returning a positive on a case where the answer is known. Ours cost
fifteen minutes of CPU time, against the four experiment generations it would
have saved.

Whether the model stores facts when the corpus is sized to its capacity rather
than twice over needs one more training run of about a day, and the corpus is
built. Until then our null marks a failed precondition rather than evidence
against the filtering argument.

\section{Limitations}
\label{sec:lim}

All model-dependent results come from one model size, one training run, and one
random seed, on synthetic data. The difficulty measurements in
Section~\ref{sec:ctrl} used an early checkpoint of the run later analysed, so
the comparison between conditions is unbiased but the absolute values are not
independent of the model. The storage result uses 200 people per group at one
checkpoint. We use \emph{stored} there in a narrow sense: that training on a
person's attribute lowers teacher-forced loss on that attribute's true value.
The measurement cannot exclude a representation that never surfaces in the
training objective, and our validation establishes sensitivity to the first
kind only. Because teacher forcing supplies the true prefix of a multi-token
value, later tokens of a value can be predicted without knowing whose value it
is, which makes the measurement conservative in a direction we have not
quantified. We claim no novelty for the
padding defect, a known failure mode of hand-written decoders that omit an
attention mask, and we have not surveyed how often the other three occur
elsewhere.

\section*{Appendix: Checks We Now Run}

None of these requires a GPU.

\small
\begin{enumerate}\itemsep0em\parskip0em\leftmargin1em
\item Decode a prompt alone and inside a batch of mixed-length prompts, and
  require identical output.
\item Save generated text, not only parsed answers. Ours stored the parse and
  discarded the text, so four rounds of invented documents left no trace.
\item Report how hard the tokens a control removes are, not only how many, and
  price each matching criterion separately.
\item Record requested data proportions beside realised ones, kept separate
  from the budget the pipeline mutates.
\item Check benchmarks for saturation as well as for floor effects.
\item Before believing a null, make the instrument produce a positive. A
  negative control alone will not detect an insensitive instrument.
\item Reconcile two measurements of the same quantity numerically before
  trusting either.
\item Give a replacement instrument the validation you demanded of the one it
  replaced.
\end{enumerate}

\begin{thebibliography}{2}
\bibitem{az} Zeyuan Allen-Zhu and Yuanzhi Li.
  Physics of Language Models: Part 3.3, Knowledge Capacity Scaling Laws.
  \emph{ICLR}, 2025. arXiv:2404.05405.
\bibitem{igsm} Tian Ye, Zicheng Xu, Yuanzhi Li, and Zeyuan Allen-Zhu.
  Physics of Language Models: Part 2.1, Grade-School Math and the Hidden
  Reasoning Process. \emph{ICLR}, 2025. arXiv:2407.20311.
\end{thebibliography}

\end{document}
